% Final section; metrics, class reports, resampling counts, and categorical
% audits were checked against the canonical M2a/M2b experiment artifacts.

\subsection{Improvement Method 2: Class Imbalance}
\label{sec:f2}

\subsubsection{Motivation}
\label{sec:f2-motivation}

The target variable of the dataset is unevenly distributed across the three outcome
classes: \textit{Graduate} is the majority class, \textit{Dropout} is the second
largest, and \textit{Enrolled} is a clear minority. On the shared held-out test set
of 885 students, the baseline model M0 correctly classifies 338 of 442
\textit{Graduate} students and 193 of 284 \textit{Dropout} students, but only 61 of
159 \textit{Enrolled} students. Consequently, although M0 achieves an overall test
accuracy of $0.668927$, its recall on the \textit{Enrolled} class is only
$0.383648$. This is a material limitation whenever the intended application requires
identifying students who are not on track to graduate on time, rather than merely
maximizing the number of correct predictions on the majority class.

Improvement Method 2 (M2) evaluates two complementary ways of increasing a decision
tree's sensitivity to minority classes: reweighting the training objective
(\textbf{M2a}) and synthetically rebalancing the training distribution
(\textbf{M2b}). Both configurations reuse the project's single stratified
train/test split
(3{,}539 training rows and 885 held-out test rows, 90 features after one-hot
encoding, \texttt{random\_state=42}); no feature scaling and no independent
re-splitting are performed in the notebook, consistent with the project-wide data
contract.

\subsubsection{Method description}
\label{sec:f2-method}

\textbf{M2a - class weighting.} M2a trains a single
\texttt{DecisionTreeClassifier} with \texttt{class\_weight=}\allowbreak\texttt{`balanced'} and
\texttt{random\_state=42}, fit directly on the original (non-resampled) training
data:
\begin{quote}
\texttt{DecisionTreeClassifier(class\_weight='balanced', random\_state=42)}
\end{quote}
Under this setting, scikit-learn scales each class's contribution to the impurity
criterion inversely to its frequency in the training data. No synthetic samples are
created; the sample size and the feature space are identical to those used by M0.

\textbf{M2b - Synthetic Minority Over-sampling Technique (SMOTE).} M2b applies
SMOTE \citep{chawla2002smote,imblearn_smote_api} to the training split only, after
the shared split has already been produced, and fits an otherwise unmodified
decision tree on the resampled data:
\begin{quote}
\texttt{X\_train\_smote, y\_train\_smote = SMOTE(} \\
\indent \texttt{sampling\_strategy="auto", random\_state=42, k\_neighbors=5} \\
\texttt{).fit\_resample(X\_train, y\_train)} \\
\texttt{DecisionTreeClassifier(random\_state=42)} \\
\indent \texttt{.fit(X\_train\_smote, y\_train\_smote)}
\end{quote}
All three SMOTE hyperparameters are declared explicitly rather than left to their
library defaults, so that the resulting artifact does not silently change if a
future \texttt{imbalanced-learn} release alters its default behavior
\citep{imblearn_oversampling_guide}. Before resampling, the training split contains
1{,}137 \textit{Dropout}, 635 \textit{Enrolled}, and 1{,}767 \textit{Graduate}
samples. After resampling, every class contains exactly 1{,}767 samples, for a total
of 5{,}301 training rows (1{,}762 synthetic rows).

\textbf{Leakage safeguard.} SMOTE is applied strictly \emph{after} the train/test
split and is fit exclusively on \texttt{X\_train}/\texttt{y\_train}. Exposing
the test set to resampling would cause data leakage, so the experiment also computes a
value-, index-, shape-, column-order-, and dtype-aware fingerprint of
\texttt{X\_train}, \texttt{X\_test}, \texttt{y\_train}, and \texttt{y\_test}
immediately before calling \texttt{SMOTE.fit\_resample()}, and asserts that all four
fingerprints are bit-for-bit unchanged immediately afterwards. The evaluation
itself is likewise performed only on the original, untouched \texttt{X\_test}/
\texttt{y\_test}; the resampled training data is used exclusively to fit the model,
never to compute any reported metric.

\subsubsection{Modified tree / model setting}
\label{sec:f2-figures}

\begin{figure}[H]
  \centering
  \includegraphics[width=0.75\textwidth]{D_cm_M2a.png}
  \caption{Confusion matrix of M2a (balanced class weights) on the 885-sample
  held-out test set.}
  \label{fig:cm-m2a}
\end{figure}

\begin{figure}[H]
  \centering
  \includegraphics[width=0.75\textwidth]{D_cm_M2b.png}
  \caption{Confusion matrix of M2b on the 885-sample held-out test set.}
  \label{fig:cm-m2b}
\end{figure}

M2b correctly classifies 201 Dropout and 74 Enrolled students, exceeding
both M0 (193 and 61) and M2a (193 and 54). It correctly classifies 334
Graduate students, more than M2a (329) but fewer than M0 (338).

\begin{figure}[H]
  \centering
  \includegraphics[width=1\textwidth]{D_tree_M2a.png}
  \caption{M2a decision tree, truncated to the first four levels for
  legibility. The fitted tree has a true depth of 28 and 696 leaves.}
  \label{fig:tree-m2a}
\end{figure}

\begin{figure}[H]
  \centering
  \includegraphics[width=1\textwidth]{D_tree_M2a_full.png}
  \caption{Full M2a decision tree (depth 28, 696 leaves).}
  \label{fig:tree-m2a-full}
\end{figure}

The full M2a view is not intended for node-by-node reading. Reweighting the
classes changes which splits the tree prioritizes, but the fitted tree remains
close to M0 in structural size (depth 27 and 634 leaves for M0).

\begin{figure}[H]
  \centering
  \includegraphics[width=1\textwidth]{D_tree_M2b.png}
  \caption{M2b decision tree displayed through the first four levels.}
  \label{fig:tree-m2b}
\end{figure}

M2b was trained on the SMOTE-resampled training set of 5{,}301 rows and
evaluated exclusively on the original held-out set. Its fitted depth is 39,
with 847 leaves; the display limit does not alter the model.

\begin{figure}[H]
  \centering
  \includegraphics[width=1\textwidth]{D_tree_M2b_full.png}
  \caption{Full M2b decision tree (depth 39, 847 leaves).}
  \label{fig:tree-m2b-full}
\end{figure}

The fitted M2b tree is larger than both M0 and M2a. SMOTE increases the
training set from 3{,}539 to 5{,}301 rows, but this experiment does not isolate
resampled row count as the sole cause of the increased structural complexity.

\subsubsection{Accuracy and error rate}
\label{sec:f2-results}

Table~\ref{tab:f2-main-metrics} reports the full set of held-out test metrics for
M0, M2a, and M2b, together with the change relative to the baseline. Table~
\ref{tab:f2-per-class} reports per-class precision and F1-score, which are not part
of the shared experiment-results artifact schema but are read directly from the two
classification reports.

\begin{table}[H]
\centering
\caption{Held-out test performance: M0 baseline versus M2a and M2b.}
\label{tab:f2-main-metrics}
\begin{tabular}{lrrrrr}
\toprule
\textbf{Metric} & \textbf{M0} & \textbf{M2a} & \textbf{M2b} & \textbf{M2a $-$ M0} & \textbf{M2b $-$ M0} \\
\midrule
Test accuracy        & 0.668927 & 0.650847 & \textbf{0.688136} & $-0.018079$ & \textbf{$+0.019209$} \\
Error rate            & 0.331073 & 0.349153 & \textbf{0.311864} & $+0.018079$ & \textbf{$-0.019209$} \\
Recall (Dropout)      & 0.679577 & 0.679577 & \textbf{0.707746} & $+0.000000$ & \textbf{$+0.028169$} \\
Recall (Enrolled)     & 0.383648 & 0.339623 & \textbf{0.465409} & $-0.044025$ & \textbf{$+0.081761$} \\
Recall (Graduate)     & \textbf{0.764706} & 0.744344 & 0.755656 & $-0.020362$ & $-0.009050$ \\
Precision (macro)     & 0.607642 & 0.584250 & \textbf{0.636513} & $-0.023392$ & \textbf{$+0.028871$} \\
Recall (macro)         & 0.609310 & 0.587848 & \textbf{0.642937} & $-0.021462$ & \textbf{$+0.033627$} \\
F1-score (macro)       & 0.608271 & 0.585409 & \textbf{0.638747} & $-0.022862$ & \textbf{$+0.030475$} \\
ROC-AUC (macro, OvR)   & 0.719456 & 0.705321 & \textbf{0.742122} & $-0.014135$ & \textbf{$+0.022667$} \\
Tree depth / leaves   & 27 / 634 & 28 / 696 & 39 / 847 & - & - \\
\bottomrule
\end{tabular}
\end{table}

\vspace{0.3cm}

\begin{table}[H]
\centering
\caption{Per-class precision and F1-score for M2a and M2b.}
\label{tab:f2-per-class}
\begin{tabular}{lrrrr}
\toprule
\textbf{Class} & \textbf{M2a Precision} & \textbf{M2b Precision} & \textbf{M2a F1} & \textbf{M2b F1} \\
\midrule
Dropout   & 0.6328 & \textbf{0.7003} & 0.6553 & \textbf{0.7040} \\
Enrolled  & 0.3234 & \textbf{0.4044} & 0.3313 & \textbf{0.4327} \\
Graduate  & 0.7966 & 0.8048 & 0.7696 & 0.7795 \\
\bottomrule
\end{tabular}
\end{table}

Notably, M2b does not merely predict \textit{Enrolled} more often at the expense of
precision: its \textit{Enrolled} precision \emph{also} increases relative to M2a
(from $0.3234$ to $0.4044$), so that when M2b predicts a student as
\textit{Enrolled} it is more likely to be correct, not merely more willing to
guess that label. This is the reason the \textit{Enrolled} F1-score of M2b
($0.4327$) is substantially higher than that of M2a ($0.3313$).

\subsubsection{Why this does (or does not) improve the result}
\label{sec:f2-discussion}

\textbf{M2a does not achieve the goal of this experiment on this split.} Test
accuracy falls by $1.81$ percentage points relative to M0; recall on
\textit{Dropout} is unchanged, recall on \textit{Enrolled} falls by $4.40$
percentage points, and recall on \textit{Graduate} falls by $2.04$ percentage
points. Because both overall accuracy and the recall of the classes this
improvement is meant to prioritize either fall or stay flat, there is no evidence
from the held-out test set that \texttt{class\_weight='balanced'} alone improves
M0 in this configuration. This negative result is reported in full rather than
omitted, since selectively reporting only favorable configurations would
misrepresent the effectiveness of class weighting for this dataset.

\textbf{On this held-out split, M2b improves the target metrics.} Recall on
\textit{Enrolled} rises from
$0.383648$ to $0.465409$ (61/159 $\rightarrow$ 74/159 students correctly
identified), and recall on \textit{Dropout} rises from $0.679577$ to $0.707746$
(193/284 $\rightarrow$ 201/284). \textit{Graduate} recall falls only marginally,
from $0.764706$ to $0.755656$. All macro-averaged metrics improve simultaneously
(macro F1: $0.608271 \rightarrow 0.638747$; macro recall:
$0.609310 \rightarrow 0.642937$). This is consistent with the intended mechanism
of SMOTE: \textit{Enrolled} has the fewest original observations in the training
data, so synthetic interpolation in its feature-space neighbourhood gives the tree
more candidate split points with which to separate \textit{Enrolled} from the
other two classes. This is a mechanistic explanation specific to this experiment,
not a causal claim that every dataset or split would show an identical gain. It is
also worth noting explicitly that M2b's tree is substantially larger than M0's
(depth 39 versus 27, 847 versus 634 leaves): the recall improvement is not the
result of a simpler, more interpretable tree, since simplification is not the
objective of this improvement method (that objective is addressed separately by
Improvement Method 1, cost-complexity pruning).

An overall test-accuracy \emph{increase} for M2b is a welcome but secondary
finding here. Because this experiment specifically targets class-wise detection,
the primary evidence is the per-class recall table above. On this held-out set,
M2b has higher Dropout and Enrolled recall than both M0 and M2a.

\subsubsection{Limitation: vanilla SMOTE on encoded categorical features}
\label{sec:f2-limitation}

The shared preprocessing pipeline one-hot encodes only four
categorical groups (\textit{Marital Status}, \textit{Application mode},
\textit{Course}, \textit{Previous qualification}); the remaining high-cardinality
categorical variables - \textit{Mother's/Father's qualification},
\textit{Mother's/Father's occupation}, \textit{Nacionality} (nominal), and
\textit{Application order} (ordinal) - are retained as integer codes. Vanilla
SMOTE treats every input feature as continuous and linearly interpolates between a
sample and its nearest neighbours, which can in principle generate values that do
not correspond to any real category.

This effect was measured directly on the 1{,}762 synthetic rows that M2b actually
uses for training, rather than assumed.

\begin{table}[H]
\centering
\caption{Out-of-vocabulary category codes among the 1{,}762 synthetic rows, for
the six integer-coded categorical columns.}
\label{tab:f2-code-audit}
\begin{tabular}{llrr}
\toprule
\textbf{Column} & \textbf{Type} & \textbf{Non-integer rate} & \textbf{Out-of-vocabulary rate} \\
\midrule
Mother's qualification & nominal & 0.00\% & 3.01\% \\
Father's qualification & nominal & 0.00\% & 2.89\% \\
Mother's occupation     & nominal & 0.00\% & 1.19\% \\
Father's occupation     & nominal & 0.00\% & 2.04\% \\
Nacionality              & nominal & 0.00\% & 2.38\% \\
Application order        & ordinal & 0.00\% & 0.00\% \\
\bottomrule
\end{tabular}
\end{table}

\vspace{0.3cm}

\begin{table}[H]
\centering\small
\caption{Interpolation outcome for the four one-hot-encoded groups, measured on
the 1{,}762 synthetic rows. Percentages are row-wise: every synthetic row falls
into exactly one of the first three columns.}
\label{tab:f2-onehot-audit}
\begin{tabular}{lrrrrr}
\toprule
\textbf{One-hot group} & \textbf{All-zero} & \textbf{Multi-hot} & \textbf{Fractional} & \textbf{Any violation} & \textbf{Valid one-hot} \\
\midrule
Marital Status          & 14.7\% & 0.0\% & 0.0\% & 14.7\% & 85.3\% \\
Application mode        & 65.4\% & 0.0\% & 0.0\% & 65.4\% & 34.6\% \\
\textbf{Course}          & \textbf{83.0\%} & 0.0\% & 0.0\% & \textbf{83.0\%} & 17.0\% \\
Previous qualification  & 24.8\% & 0.0\% & 0.0\% & 24.8\% & 75.2\% \\
\bottomrule
\end{tabular}
\end{table}

The two measurements must be read differently.

For the \textbf{six integer-coded columns}, the audit detects no non-integer
values whatsoever in the synthetic data. This result must \emph{not} be
interpreted as "SMOTE does not affect these columns": it is very plausible that
interpolated fractional values were truncated to integers during
\texttt{fit\_resample()}'s dtype-restoration step, which would make this
measurement, at the level of raw values, unable to detect the interpolation that
still occurred. Even where a value is an integer, it can still fail to correspond
to any real category - and the out-of-vocabulary-rate column above shows exactly
this: up to $3.01\%$ of synthetic values for \textit{Mother's qualification} are
integer codes that never appear anywhere in the real training data.

For the \textbf{four one-hot groups}, the measurement is unambiguous and fully
quantitative. \texttt{imbalanced-learn} 0.14.2's internal \texttt{ArraysTransformer}
restores the output DataFrame to the exact dtype of its input columns via
\texttt{astype()} \citep{imblearn_arrays_transformer}. Because the dummy columns
produced have integer dtype, any fractional value produced
by SMOTE's linear interpolation between two different one-hot vectors is truncated
toward zero rather than rounded. Given that a convex combination of two valid
one-hot vectors can never itself exceed 1 in any single component, this truncation
can only push a component down to $0$, never up past $1$: the audit accordingly
finds a fractional-component rate of $0.0\%$ and a multi-hot rate of $0.0\%$ in
every group, and \emph{100\% of the observed violations are all-zero vectors}
(no category active) rather than any mixture of multiple categories. \textit{Course}'s high cardinality is a plausible contributor to its high all-zero violation rate ($83.0\%$) under this encoding-and-restoration path. The audit does not isolate category count as the sole cause.

This is a documented limitation of vanilla SMOTE when applied to mixed
continuous/categorical representations; the official \texttt{imbalanced-learn}
documentation recommends \texttt{SMOTENC} for exactly this situation
\citep{imblearn_oversampling_guide}. Vanilla SMOTE was retained as one of the group's selected class-imbalance experiments. The assignment brief suggests handling class imbalance as a possible improvement direction but does not mandate SMOTE specifically. The M2b results reported above should
therefore be read together with this limitation: the recall and precision
improvements are measured on the real, untouched held-out test set, but
part of the synthetic training data the tree learns from does not represent a
"valid" student under the true meaning of the affected categorical columns.

\subsubsection{Accuracy-recall trade-off}
\label{sec:f2-tradeoff}

In a class-imbalanced classification problem, overall accuracy alone cannot be the
sole acceptance criterion. A model may legitimately trade a reduction in overall
accuracy for less majority-class dominance; when this happens,
per-class recall - not a return to M0's accuracy figure - is the correct measure
of success.

M2a illustrates that this trade-off is not automatic: its accuracy falls
\emph{and} its \textit{Enrolled} recall falls, so the result is reported plainly
rather than selectively. M2b, by contrast, is not a trade-off at all on this
held-out split: it simultaneously increases test accuracy ($0.668927 \rightarrow 0.688136$), \textit{Dropout}
recall, and \textit{Enrolled} recall, at the cost of a $0.90$ percentage-point
reduction in \textit{Graduate} recall. If the evaluation priority is detecting
students who are at risk of dropping out or remaining indefinitely enrolled, M2b
is therefore the preferable configuration among M0, M2a, and M2b on this held-out
split.

These findings are drawn from a single fixed held-out split, and the synthetic
portion of M2b's training data carries the categorical-representation limitation
discussed above. Before any real-world deployment, results should be validated
across multiple splits or cohorts, \texttt{SMOTENC} should be considered as an
additional experiment, and performance disparities across the dataset's
sensitive groups should be examined. Within the scope of this experiment and
the project's shared data split, however, M2b answers the improvement question
affirmatively: the
minority \textit{Enrolled} class and the \textit{Dropout} class are both
recognized more reliably when SMOTE is applied strictly to the training split.

\subsubsection{Conclusion}
\label{sec:f2-conclusion}

Balanced class weighting (M2a) does not improve upon the M0 baseline on this
held-out test set: overall accuracy, \textit{Enrolled} recall, and
\textit{Graduate} recall all fall, while \textit{Dropout} recall is unchanged.
Applying SMOTE to the training split only (M2b) increases \textit{Enrolled}
recall by $8.18$ percentage points and \textit{Dropout} recall by $2.82$
percentage points relative to M0, while simultaneously raising \textit{Enrolled}
precision from $0.3567$ to $0.4044$ and lowering the overall error rate from
$0.331073$ to $0.311864$. M2b is therefore adopted as the more successful
class-imbalance configuration within the scope of this experiment. M2a is
retained in this report rather than discarded, both for transparency - not every
class-balancing technique improves a given model on a given dataset - and to
document, alongside M2b, the quantitative categorical-representation limitation
of vanilla SMOTE rather than presenting only favorable figures.
